% from ega0-16-fr.tex
\oldpage[0\textsubscript{IV-1}]{24}
Si $\mathfrak{n}$ est le noyau de $\rho$, $\rho(A)$ est isomorphe à
$A/\mathfrak{n}$, donc $\dim\rho(A)\leq\dim(A)$ par (16.1.2.1), et comme $B$
est entier sur $\rho(A)$, on peut se borner à démontrer la seconde assertion
lorsque $A\subset B$. Si
$\mathfrak{p}'_0\subset\mathfrak{p}'_1\subset\cdots\subset\mathfrak{p}'_n$
est une chaîne d'idéaux premiers de $B$, les $\mathfrak{p}'_i\cap A$ sont
alors des idéaux premiers distincts dans $A$ (Bourbaki, \emph{Alg.~comm.},
chap.~V, \S~2, no~1, cor.~1 de la prop.~1), donc forment une chaîne d'idéaux
premiers. Inversement, pour toute chaîne
$\mathfrak{p}_0\subset\mathfrak{p}_1\subset\cdots\subset\mathfrak{p}_m$
d'idéaux premiers de $A$, on sait qu'il existe pour chaque $i$ un idéal
premier $\mathfrak{p}'_i$ de $B$ tel que
$\mathfrak{p}'_i\cap A=\mathfrak{p}_i$ et que
$\mathfrak{p}'_i\subset\mathfrak{p}'_{i+1}$ pour $0\leq i\leq m-1$
(\emph{loc.~cit.}, cor.~2 du th.~1). D'où la conclusion.
\end{proof}

\begin{proposition}[16.1.6]
\label{0.16.1.6-fr}
Soient $B$ un anneau intègre, $A$ un sous-anneau local de $B$,
$\mathfrak{m}$ son idéal maximal. Supposons que $A$ soit intégralement clos et
que $B$ soit entier sur $A$~; alors, pour tout idéal maximal $\mathfrak{n}$ de
$B$, on a $\dim(B_{\mathfrak{n}})=\dim(A)$.
\end{proposition}

\begin{proof}
La première partie du raisonnement de (16.1.5) montre que
$\dim(B_{\mathfrak{n}})\leq\dim(A)$. Inversement, considérons une chaîne
d'idéaux premiers
$\mathfrak{p}_0\subset\mathfrak{p}_1\subset\cdots\subset
\mathfrak{p}_k=\mathfrak{m}$ de $A$~; on sait que
$\mathfrak{n}\cap A=\mathfrak{m}$ (Bourbaki, \emph{Alg.~comm.}, chap.~V,
\S~2, no~1, prop.~1) et en vertu des hypothèses faites, on peut appliquer le
second théorème de Cohen-Seidenberg (Bourbaki, \emph{Alg.~comm.}, chap.~V,
\S~2, no~4, th.~3), qui prouve l'existence d'une chaîne d'idéaux premiers
$\mathfrak{p}'_0\subset\mathfrak{p}'_1\subset\cdots\subset
\mathfrak{p}'_k=\mathfrak{n}$ de $B$ telle que
$\mathfrak{p}'_i\cap A=\mathfrak{p}_i$ pour tout $i$~; on en déduit aussitôt
$\dim(B_{\mathfrak{n}})\geq\dim(A)$, d'où la proposition.
\end{proof}

\begin{env}[16.1.7]
\label{0.16.1.7-fr}
Soit $M$ un $A$-module~; on appelle dimension de $M$ et on note
$\dim_A(M)$ ou $\dim(M)$ la dimension de l'anneau $A/\operatorname{Ann}(M)$
(isomorphe à l'anneau des homothéties $A_M$ de $M$). On a donc
$\dim(M^k)=\dim(M)$ pour tout $k>0$. Si $M$ est de type fini, les idéaux
premiers de $A$ contenant $\operatorname{Ann}(M)$ sont exactement ceux qui
appartiennent à $\operatorname{Supp}(M)$ ($0_{\mathrm I}$, 1.7.4) et ce
dernier est fermé dans $\operatorname{Spec}(A)$, donc on a alors
\[
\label{0.16.1.7.1-fr}
  \dim(M)=\dim(\operatorname{Supp}(M))\leq\dim(A).
\tag{16.1.7.1}
\]
Si $M$, $N$ sont deux $A$-modules tels que $N\subset M$, on a
\[
\label{0.16.1.7.2-fr}
  \dim(N)\leq\dim(M),\qquad \dim(M/N)\leq\dim(M).
\tag{16.1.7.2}
\]
Si $S$ est une partie multiplicative de $A$ et si $M$ est de type fini, on a
\[
\label{0.16.1.7.3-fr}
  \dim_{S^{-1}A}(S^{-1}M)\leq\dim_A(M).
\tag{16.1.7.3}
\]
En effet, on a
$\operatorname{Ann}(S^{-1}M)=S^{-1}\operatorname{Ann}(M)$ (Bourbaki,
\emph{Alg.~comm.}, chap.~II, \S~2, no~4) et
$S^{-1}A/S^{-1}\operatorname{Ann}(M)=S^{-1}(A/\operatorname{Ann}(M))$.

Par ailleurs, toute chaîne maximale d'idéaux premiers contenant
$\operatorname{Ann}(M)$ se termine par un idéal maximal $\mathfrak{m}$, donc
sa longueur est égale à celle de la chaîne des idéaux premiers correspondants
de $A_{\mathfrak{m}}$ (qui contiennent
$\operatorname{Ann}(M_{\mathfrak{m}})=(\operatorname{Ann}(M))_{\mathfrak{m}}$)~;
de cette remarque et de (16.1.7.3), on tire
\[
\label{0.16.1.7.4-fr}
  \dim_A(M)=\sup_{\mathfrak{m}}
  (\dim_{A_{\mathfrak{m}}}(M_{\mathfrak{m}}))
\tag{16.1.7.4}
\]
où $\mathfrak{m}$ parcourt l'ensemble des idéaux maximaux de $A$.

On dit que $M$ est caténaire (resp. équidimensionnel, équicodimensionnel,
biéquidimensionnel) lorsque $A/\operatorname{Ann}(M)$ l'est.
\end{env}

\begin{proposition}[16.1.8]
\label{0.16.1.8-fr}
